%%%%%%%%%%  scan idx 353  |  folio 340  %%%%%%%%%%
\(B\) qui
envoie \(N\) sur \(N''\), et \(x\) sur \(x''\)~; d'après la remarque précédente, l'élément obtenu ne dépend pas de l'entier \(d\) tel que \(\lambda^k(N'')=0\) pour \(k>d\). En utilisant le même homomorphisme, on voit que la relation (5.3.1) est vérifiée par \(N''\) et \(x''\).

\textbf{5.4.} Considérons le \(\lambda\)-anneau engendré sur \(\mathbb Z\) par les indéterminées \(N,x,y\), soumises aux relations \(\lambda^k(N)=0\) pour \(k>d\). De (5.3.1) résulte :
\begin{align*}
  \lambda^p(N,x+y)\cdot\lambda_{-1}(N)
  &=\lambda^p\bigl((x+y)\cdot\lambda_{-1}(N)\bigr)\\
  &=\sum_{i+j=p}\lambda^i(x\cdot\lambda_{-1}(N))\cdot\lambda^j(y\cdot\lambda_{-1}(N))\\
  &=\sum_{\substack{i+j=p\\ i,j\ne0}}\lambda^i(N,x)\cdot\lambda^j(N,y)\cdot(\lambda_{-1}(N))^2
    +(\lambda^p(N,x)+\lambda^p(N,y))\lambda_{-1}N
\end{align*}

Comme l'anneau considéré est intègre, on en désuit la relation suivante, valable par substitution dans tout \(\lambda\)-anneau, et pour tout élément \(N\) tel que \(\lambda^k(N)=0\) pour \(k>d\) :
\begin{equation}\tag{5.4.1}
  \lambda^p(N,x+y)=\sum_{\substack{i+j=p\\ i,j\ne0}}\lambda^i(N,x)\cdot\lambda^j(N,y)\cdot\lambda_{-1}(N)+\lambda^p(N,x)+\lambda^p(N,y).
\end{equation}

Considérons maintenant le \(\lambda\)-anneau engendré sur \(\mathbb Z\) par des indéterminées \(N,N',x\) soumises aux relations \(\lambda^k(N)=0\) pour \(k>d\), et \(\lambda^k(N')=0\) pour \(k>d'\). On a alors \(\lambda^k(N+N')=0\) pour \(k>d+d'\)~; on peut donc définir les \(\lambda^p(N+N',x)\). On obtient :
\[
  \lambda^p(N+N',x)\cdot\lambda_{-1}(N+N')=\lambda^p(x\cdot\lambda_{-1}(N+N')).
\]
Par additivité de \(\lambda_t\), on a : \(\lambda_{-1}(N+N')=\lambda_{-1}(N)\cdot\lambda_{-1}(N')\). On en déduit :
\begin{align*}
  \lambda^p(N+N',x)\cdot\lambda_{-1}(N)\cdot\lambda_{-1}(N')
  &=\lambda^p(x\cdot\lambda_{-1}(N)\cdot\lambda_{-1}(N'))\\
  &=\lambda^p(N',x\cdot\lambda_{-1}(N))\cdot\lambda_{-1}(N').
\end{align*}
